
\section{\textbf{\texttt{filtered()}} ORM Method}

\begin{frame}[fragile]{\textbf{\texttt{filtered()}} Method}
	\begin{itemize}
		\item The \textbf{\texttt{filtered()}} method filters an existing recordset in Python.
		\item Unlike \texttt{search()}, it does \textbf{not} run a new SQL query by domain.
		\item It returns a new recordset containing only records that match a condition.
\begin{block}{Method Syntax}
\begin{lstlisting}
def filtered(self, func):
	return super().filtered(func)
\end{lstlisting}
	\begin{itemize}
		\item Important parts: \textbf{\texttt{self}}, \textbf{\texttt{func}}.
	\end{itemize}
\end{block}
	\end{itemize}
\end{frame}

\begin{frame}[fragile]{Breaking Down the Method Signature}
	\begin{block}{}
		\begin{itemize}
			\item \textbf{\texttt{self}}: The starting recordset.
			\item \textbf{\texttt{func}}: Either
			\begin{itemize}
				\item a function / lambda that returns True/False, or
				\item a field name string (truthy check).
			\end{itemize}
\begin{lstlisting}
students.filtered(lambda s: s.age >= 18)
students.filtered('email')   # keep records with a truthy email
\end{lstlisting}
		\end{itemize}
	\end{block}
\end{frame}

\begin{frame}[fragile]{What Does \texttt{filtered()} Return?}
	\begin{block}{}
		\begin{itemize}
			\item It returns a recordset.
\begin{lstlisting}
students = self.env['student.student'].search([])
adults = students.filtered(lambda s: s.age >= 18)
print(adults)
\end{lstlisting}
			Output is (example):
\begin{lstlisting}
student.student(15, 22)
\end{lstlisting}
		\end{itemize}
	\end{block}
\end{frame}

\begin{frame}[fragile]{Filtered with a Lambda}
\begin{lstlisting}
students = self.env['student.student'].browse([15, 16, 17])

males = students.filtered(lambda s: s.gender == 'male')
with_email = students.filtered(lambda s: bool(s.email))
active_adults = students.filtered(
	lambda s: s.active and s.age >= 18
)
\end{lstlisting}
	\begin{itemize}
		\item The lambda receives one record at a time.
		\item Return \texttt{True} to keep the record.
	\end{itemize}
\end{frame}

\begin{frame}[fragile]{\texttt{filtered()} vs \texttt{search()}}
	\begin{itemize}
		\item Use \texttt{search()} when filtering from the whole table with a domain (SQL).
		\item Use \texttt{filtered()} when you already have a recordset and need a local filter.
\begin{lstlisting}
# From database
adults = self.env['student.student'].search([('age', '>=', 18)])

# From an existing recordset
adults = students.filtered(lambda s: s.age >= 18)
\end{lstlisting}
	\end{itemize}
\end{frame}

\begin{frame}[fragile]{Method Call}
\begin{block}{Method Call}
\begin{lstlisting}
adults = students.filtered(lambda student: student.age >= 18)
\end{lstlisting}
	\begin{itemize}
		\item Here, \textbf{\texttt{.filtered(...)}} is the \textbf{method call}.
		\item Return type reminder: \textbf{recordset}.
	\end{itemize}
\end{block}
\end{frame}
